\documentclass[10pt,aspectratio=169]{beamer}
\usepackage{bq-theme}
\usepackage{pgfplots}
\pgfplotsset{compat=1.18}
\hypersetup{pdftitle={Local Responses and Optimized Values}}
\newcommand{\dd}{\mathrm{d}}

% Replace this small illustrative example with the supplied research material.
% The shared source deliberately contains no McGill findings or calibrations.
\newcommand{\SupplyFigure}[1]{%
\begin{tikzpicture}
\begin{axis}[width=.82\textwidth,height=.54\textheight,
  xmin=0,xmax=1.7,ymin=0,ymax=1.7,axis lines=left,
  xlabel={Quantity $q$},ylabel={Marginal cost / price},
  xtick={0,1,1.3},ytick={0,1,1.3},tick label style={font=\small},
  label style={font=\small},clip=false]
  \addplot[black,thick,domain=0:1.6,samples=2]{x};
  \node[anchor=west,font=\small] at (axis cs:1.2,1.55) {$C'(q)=q$};
  \addplot[BQBlue,dashed] coordinates {(0,1) (1,1) (1,0)};
  \addplot[BQBlue,only marks,mark=*] coordinates {(1,1)};
  \node[anchor=south east,font=\small,text=BQBlue] at (axis cs:1,1) {Initial optimum};
  \ifnum#1=2
    \addplot[BQRed,dashed] coordinates {(0,1.3) (1.3,1.3) (1.3,0)};
    \addplot[BQRed,only marks,mark=*] coordinates {(1.3,1.3)};
    \node[anchor=west,font=\small,text=BQRed] at (axis cs:1.32,1.27) {After $p$ rises};
  \fi
\end{axis}
\end{tikzpicture}}

\begin{document}
\begin{frame}[plain,noframenumbering,label=title]
\vspace{8mm}
\begin{center}
{\Large\color{BQBlue}Local Responses and Optimized Values\par}
\vspace{10mm}
\begin{tabular}{ccc}
Author One & Author Two & Author Three\\[2pt]
{\footnotesize Affiliation One} & {\footnotesize Affiliation Two} & {\footnotesize Affiliation Three}
\end{tabular}\par
\vspace{8mm}
{\small Research Seminar\\[3pt]Month Year}
\end{center}
\vfill
{\tiny Illustrative starter, not empirical research results.}
\end{frame}

\begin{frame}[label=main_question]{How Does a Price Change Affect Profits?}
\begin{BQList}
\item<1-> A price change has two apparent effects.
\begin{BQSub}
\item Revenue changes on the initial quantity.
\item The firm chooses a new quantity.
\end{BQSub}
\item<2-> Do we need to predict both to value a small price change?
\item<3-> Optimization makes the \textbf{initial quantity} informative.
\end{BQList}
\end{frame}

\BQAgenda{1}{A local mechanism,What the result identifies,From theory to measurement}

\begin{frame}[label=main_setup]{Begin with the Firm's Choice}
\begin{BQCompact}
\item<1-> A price-taking firm chooses output $q$ at price $p$:
\[G(p)=\max_{q\geq0}\{pq-C(q)\}.\]
\item<2-> Hold the cost function fixed and consider a differentiable interior optimum.
\begin{BQSub}
\item The first-order condition is $C'(q)=p$.
\end{BQSub}
\item<3-> We want the change in optimized profit, $\dd G$.
\end{BQCompact}
\BQBottomLinks{\backupbtn{backup_derivation}{Derivation}}
\end{frame}

\begin{frame}[label=main_result]{The Initial Quantity Values a Small Price Change}
\begin{BQCompact}[12pt]
\item<1-> Differentiate revenue and cost along the firm's choice:
\[\dd G=\textcolor{BQBlue}{q\,\dd p}+[p-C'(q)]\dd q.\]
\item<2-> At the optimum, the \textbf{quantity adjustment cancels}:
\[p-C'(q)=0\quad\Longrightarrow\quad\dd G=\textcolor{BQBlue}{q\,\dd p}.\]
\item<3-> The \textcolor{BQBlue}{gain on the initial quantity} is the first-order profit change.
\begin{BQSub}
\item No separate prediction of $\dd q$ is needed for this local valuation.
\end{BQSub}
\end{BQCompact}
\BQSource{Price-taking, fixed differentiable costs, and an interior optimum. This is a local statement.}
\BQBottomLinks{\backupbtn{backup_derivation}{Cancellation}\quad\backupbtn{backup_finite}{Finite changes}}
\end{frame}

\begin{frame}[t,label=main_figure]{Output Adjusts, but Its First-Order Profit Effect Cancels}
\BQLead{Illustrate with $C(q)=q^2/2$, so the firm chooses $q=p$.}
\begin{overlayarea}{\textwidth}{.58\textheight}
\centering
\only<1|handout:0>{\SupplyFigure{1}}
\only<2-|handout:1>{\SupplyFigure{2}}
\end{overlayarea}
\begin{itemize}
\item<3-> A small price rise values each unit initially sold at the higher price.
\end{itemize}
\BQSource{Illustrative quadratic cost example; axis values are not estimates.}
\end{frame}

\begin{frame}[label=main_limits]{Local Valuation and Finite Changes Are Different Tasks}
\begin{BQList}
\item<1-> Initial quantity identifies the \textbf{local slope} of optimized profit.
\item<2-> A finite change also depends on how output varies along the price path:
\[G(p_1)-G(p_0)=\int_{p_0}^{p_1}q(p)\,\dd p.\]
\item<3-> An empirical application must distinguish an observed quantity from an assumed response curve.
\end{BQList}
\BQBottomLinks{\backupbtn{backup_finite}{Quadratic example}}
\end{frame}

\begin{frame}[label=bq_main_end]{Conclusion}
\begin{BQList}
\item<1-> Optimization cancels the first-order effect of the chosen quantity adjustment.
\item<2-> Initial quantity values a small price change.
\item<3-> Finite-change predictions require additional information about the response path.
\end{BQList}
\end{frame}

\BQAppendix
\setappendixreturn{main_result}{Back to the envelope}
\begin{frame}[noframenumbering,label=backup_derivation]{Why Does the Quantity Adjustment Cancel?}
\begin{BQCompact}
\item Write optimized profit as $G(p)=p\,q(p)-C(q(p))$.
\item Apply the chain rule:
\[G'(p)=q(p)+p\,q'(p)-C'(q(p))q'(p).\]
\item The first-order condition $C'(q(p))=p$ leaves $G'(p)=q(p)$.
\item The calculation assumes a smooth interior solution; changing costs would add another term.
\end{BQCompact}
\end{frame}

\setappendixreturn{main_limits}{Back to local versus finite}
\begin{frame}[noframenumbering,label=backup_finite]{The Finite-Change Remainder Is Second Order}
\begin{BQCompact}
\item In the quadratic example, $q(p)=p$ and $G(p)=p^2/2$.
\item For price change $h$:
\[G(p+h)-G(p)=p\,h+\frac{h^2}{2}.\]
\item The envelope gives the linear term $q(p)h=p\,h$.
\item The remainder depends on curvature; the local formula is not an exact finite-change identity.
\end{BQCompact}
\end{frame}
\end{document}
